\section{Introduction}
\label{sec:intro}

Road traffic injury kills approximately 1.19 million people a year by the World
Health Organization's count~\cite{who2023}. Many of those deaths are
foreseeable in the seconds before they occur, which is why systems that
anticipate a collision --- rather than detect one after the fact --- are
pursued as a countermeasure, and why a growing literature reports how early
such systems can raise an alarm~\cite{zhang2025review}. Whether a reported
anticipation result establishes that a system could actually warn a driver is a
question about measurement, and it is the question this paper addresses.

Anticipation is increasingly scored by a composite metric: a weighted sum of
discrimination terms --- average precision (AP) and area under the ROC curve
(AUC) --- and an earliness term, the interval between alarm and collision read
at a fixed risk threshold. The composition is intuitive. An anticipator should
separate accident clips from ordinary driving and should raise the alarm early,
and summing terms that measure each seems a reasonable way to ask for both.

It is not, and the reason is a difficulty of type rather than of tuning. AP and
AUC are bounded on the unit interval; an earliness term counted in frames is
bounded only by the length of the clip. Summing a bounded quantity with an
unbounded one lets the unbounded one decide the ranking, and makes that ranking
depend on how the corpus was windowed. Worse, an earliness term read at a
threshold has a trivial global maximiser: a risk score that exceeds the
threshold at frame~0 and never falls attains, on every clip, the largest value
that clip admits. No method that reads the video can exceed it. A method can
only match it, and then compete on the bounded half.

That argument follows from the definitions alone. What has been missing is a
measurement of what it costs in practice, and the measurement is awkward to
obtain: the benchmarks using these metrics withhold their labels, and the
competitions using them do not publish the weights.

This paper supplies it twice over. On a live accident-anticipation benchmark
--- 1{,}417 dashcam clips, thirteen teams, a scorer we never saw --- a constant
risk score of $0.51$ that reads no pixels scores $2.35291$ against a winning
$2.50165$, matching the eighth-placed team at the precision the leaderboard
displays and exceeding five of thirteen entries. A family of submissions that
likewise read nothing, designed so that the undisclosed terms cancel when two
are differenced, then recovers the metric's parts from outside the competition:
the timing terms are worth \textbf{5.70 times} what discrimination contributes
to such a submission, and the score falls at almost exactly one sixtieth of a
point for every frame of delay.

Because that benchmark publishes no labels, it can show the metric being
exploited but cannot show anything being fixed. We therefore repeat the
measurement on an independent, fully annotated corpus of 1{,}500 clips, half of
which contain no collision. There the pathology reappears --- a constant takes
99.9\% of the available earliness at exactly chance discrimination --- and, for
the first time, the cost of the behaviour it rewards becomes measurable: a
warning that maximises this metric alarms for \textbf{60 minutes in every hour}
of ordinary driving.

The safety consequence is therefore not that the metric is untidy. The
submission maximising its earliness half is a warning that never turns off.
Whether a driver comes to trust a collision warning or to ignore it depends on
when it arrives relative to what the driver expects~\cite{abe2006alarm}, and a
warning that never stops carries no timing information at all. The metric
awards that behaviour its maximum and cannot do otherwise, because its timing
terms are defined only on clips where an accident occurs, so the burden such a
system imposes on ordinary driving falls entirely on the discrimination half.

We came to the question the awkward way: we built a training-free ensemble,
entered it in that competition, and placed ninth of thirteen. The paper is not
about that system --- it scores below the constant --- but it is what made the
metric worth examining, and it appears here as the case study.

\textbf{Contributions.} We report a measured demonstration on a live
benchmark: a video-blind constant matches the eighth-placed private score and
exceeds five of thirteen teams, our own included. We give a black-box
decomposition, in which video-blind probes built so that unknown terms cancel
under differencing separate a metric whose weights are undisclosed and whose
labels are withheld, recovering the discrimination floor, both timing terms,
the slope in the crossing frame and the accident-onset distribution. We show
that the published definition does not reproduce the scorer: four curves
crossing at the same frame, none dipping below it afterwards, score $0.091$
apart, where the stated formula requires equality. We then validate the
corrected protocol where ground truth exists; its predicted behaviour holds,
but it is beaten by a control reading only the video file's header, so we
replace its clip-level conditioning with a frame-level alarm statistic and show
that this orders submissions by what they can see. Every leaderboard score was
computed by the competition's scoring service against ground truth we do not
have; derived quantities are our arithmetic on those scores and are marked as
such.

One point of scope, because this paper prints a number beside other teams'
numbers. We make no claim about any other submission or published method: we
did not re-run anyone's system, and could not have evaluated it if we had,
because the labels are withheld from every entrant. Where a leaderboard score
is numerically identical to a control's, we state that arithmetic fact and draw
one inference --- the metric assigns the two submissions the same value. A
whole family of curves receives it, which is precisely the property under
examination.
